% !TEX root = ../00_Main/Main.tex
\documentclass[../00_Main/Main.tex]{subfiles}
\begin{document}

\label{sec:conclusion}

\subsection{Respuesta a las preguntas de investigaci\'on}

\textbf{Mejor modelo individual.} Los resultados respaldan la hip\'otesis
H1. El Transformer causal obtuvo el mayor desempe\~no medio entre los
modelos individuales, en la referencia y en los escenarios de
generalizaci\'on; en estos \'ultimos super\'o en un $3\%$ a DQN y en un
$5\%$ al DQN recurrente, diferencias de tama\~no medio. Mantuvo el
desempe\~no con severidades fuera de rango, pero no fue el m\'as estable
en su peor escenario: con secuencias largas perdi\'o un $7\%$ respecto de
su referencia, frente al $6\%$ de DQN.

\textbf{Ensambles frente al mejor modelo individual.} Los resultados no
respaldan la hip\'otesis H2. S\'olo el ensamble ponderado super\'o al Transformer,
en un $1\%$ en generalizaci\'on, una diferencia significativa pero de
tama\~no despreciable. La
compuerta del Transformer alcanz\'o un desempe\~no equivalente y los
otros cuatro ensambles quedaron por debajo. La ventaja del ensamble ponderado no se explica por
la ponderaci\'on por confianza: su acci\'on final difiere de la del
Transformer en el $61\%$ de las decisiones de la referencia, en buena
medida por la temperatura alta y el \emph{prior} de severidad.

\textbf{Ponderaci\'on por confianza.} El dise\~no no permite determinar si
ponderar por la confianza mejora a una votaci\'on simple, porque todos los
ensambles aplican esa ponderaci\'on. Los que s\'olo promedian o votan con
esos pesos obtuvieron un desempe\~no inferior al del Transformer
individual.

\subsection{Limitaciones del estudio}\label{sec:limitations}

\begin{itemize}
    \item Cada escenario se gener\'o con una \'unica semilla ($42$), lo que
    impide estimar la variabilidad de las medias y construir intervalos de
    confianza para el ordenamiento.
    \item Los hiperpar\'ametros se ajustaron sobre las mismas $64$
    secuencias de la referencia, por lo que el desempe\~no en esa
    condici\'on puede estar sobreestimado.
    \item Las arquitecturas del DQN recurrente y del Transformer se
    eligieron \emph{ad hoc}; otras configuraciones podr\'ian modificar el
    ordenamiento.
    \item La confianza se calcula sobre valores $Q$ transformados y no
    sobre probabilidades calibradas, y no se verific\'o su relaci\'on con
    el acierto de las decisiones.
\end{itemize}

\subsection{Trabajo futuro}\label{sec:future}

\begin{itemize}
    \item Repetir la evaluaci\'on con semillas independientes para
    obtener intervalos de confianza del ordenamiento de la
    Tabla~\ref{tab:global-mean}.
    \item Relacionar la confianza con el acierto de cada decisi\'on
    mediante curvas de calibraci\'on.
    \item Aprender pesos $w_k(s)$ dependientes del estado.
\end{itemize}

\subsection{Conclusi\'on general}

\enlargethispage{2\baselineskip} % evita que las \'ultimas l\'ineas queden solas en una p\'agina
En el entorno controlado de mPES, el Transformer causal es el modelo
individual con mejor desempe\~no, en las condiciones de entrenamiento y
fuera de ellas, aunque no es el que menos se degrada en su peor
escenario. La combinaci\'on de modelos lo super\'o en un solo caso y por un
margen despreciable, mediante una regla que suaviza las preferencias de
los miembros y las mezcla con un \emph{prior} de severidad. En
consecuencia, H1 no se refuta y H2 se refuta. Estas conclusiones se limitan al entorno y a los modelos
evaluados.

\biblio % Needed for referencing to working when compiling individual subfiles - Do not remove
\end{document}
